%%%%%%%%%%%%%%%%%%%%%%%%%%%%%%%%%%%%%%%%%%%%%%%%%%%%%%%%%%%%%%%%%%%%%%%%%%%%%%%%%%
\begin{frame}[fragile]\frametitle{}

\begin{center}
{\Large Google NotebookLM}
\end{center}
\end{frame}

%%%%%%%%%%%%%%%%%%%%%%%%%%%%%%%%%%%%%%%%%%%%%%%%%%%%%%%%%%%%%%%%%%%%%%%%%%%%%%%%%%
\begin{frame}[fragile]\frametitle{What is NotebookLM?}
  \begin{itemize}
    \item \textbf{One-Line Definition:} A source-grounded AI research assistant by Google Labs — it answers questions, synthesizes insights, and generates multimodal outputs \textit{exclusively from documents you provide}.
    \item \textbf{Origin Story:} First introduced in May 2023 as \textit{Project Tailwind} at Google I/O. Rebranded as NotebookLM in 2024 and rolled out broadly to researchers, students, and enterprise users.
    \item \textbf{The Core Idea:} Unlike general chatbots that draw on broad training data, NotebookLM is grounded in \textit{your} sources — minimizing hallucinations and providing inline citations linking directly to source passages.
    \item \textbf{Underlying Model:} Powered by Google Gemini (upgraded to Gemini 2.0 Flash in Dec 2024; Gemini 3 in 2025) — with higher model tiers available for paid plans.
    \item \textbf{Availability:} Web app at \href{https://notebooklm.google.com}{\texttt{notebooklm.google.com}} + mobile app · 180+ regions · 50+ languages · Free to use with a Google account.
    \item \textbf{Scale:} Globally, millions of people and tens of thousands of organizations use NotebookLM; users have generated more than 350 years' worth of Audio Overviews.
  \end{itemize}
\end{frame}

%%%%%%%%%%%%%%%%%%%%%%%%%%%%%%%%%%%%%%%%%%%%%%%%%%%%%%%%%%%%%%%%%%%%%%%%%%%%%%%%%%
\begin{frame}[fragile]\frametitle{The Three-Panel Interface}
  \begin{itemize}
    \item \textbf{Redesigned Dec 2024:} NotebookLM introduced a three-column layout that organizes the entire research workflow into three focused areas.
    \item \textbf{Sources Panel (Left):} The central knowledge base for your project. Add, manage, and browse all uploaded sources. Click any source to open the full document viewer. Select specific sources to scope chat responses.
    \item \textbf{Chat Panel (Center):} Conversational AI interface grounded in your sources. Ask questions, request summaries, or generate custom content. Every response includes clickable inline citations that jump directly to the relevant passage in the original document.
    \item \textbf{Studio Panel (Right):} One-click content generation from your sources — Audio Overviews, Video Overviews, Mind Maps, Reports, Flashcards, Quizzes, Slide Decks, Infographics, and Study Guides. All outputs are saved in the Studio for ongoing access.
    \item \textbf{Adaptive Layout:} Expand the source viewer and notes editor side by side for deep reading, or collapse panels to focus on chat while an Audio Overview plays in the background.
  \end{itemize}
\end{frame}

%%%%%%%%%%%%%%%%%%%%%%%%%%%%%%%%%%%%%%%%%%%%%%%%%%%%%%%%%%%%%%%%%%%%%%%%%%%%%%%%%%
\begin{frame}[fragile]\frametitle{Supported Source Types}
  \begin{itemize}
    \item \textbf{Documents:} PDFs (up to 200MB / 500,000 words each), Google Docs, Google Slides, Google Sheets, Microsoft Word (\texttt{.docx}), plain text files. Note: copy-protected PDFs will fail to import.
    \item \textbf{Web \& Video:} Website URLs (public pages; paywalled or scraping-blocked sites will fail), YouTube video links (NotebookLM extracts the transcript automatically).
    \item \textbf{Audio:} MP3, WAV, AAC, OGG, OPUS, and 20+ other formats containing clear speech (music-only files are not supported).
    \item \textbf{Images:} Upload images as sources for visual content analysis.
    \item \textbf{Notes:} Manually authored notes within a notebook can be converted into sources, letting you build on AI analysis with your own insights.
    \item \textbf{Google Drive Integration:} Pull in Google Docs, Slides, and Sheets directly from Drive — no manual download required.
    \item \textbf{Limitations:} Excel files are not directly supported (export as PDF as a workaround). Sources are static — updated versions must be manually replaced.
  \end{itemize}
\end{frame}

%%%%%%%%%%%%%%%%%%%%%%%%%%%%%%%%%%%%%%%%%%%%%%%%%%%%%%%%%%%%%%%%%%%%%%%%%%%%%%%%%%
\begin{frame}[fragile]\frametitle{Audio Overviews — The Viral Feature}
  \begin{itemize}
    \item \textbf{What It Is:} Upload any set of documents and NotebookLM generates a podcast-style conversation between two lifelike AI hosts who summarize, debate, and contextualize your content.
    \item \textbf{Launched Sep 2024:} Immediately viral — Spotify used it for personalized 2024 Wrapped podcast overviews. Users generated over 350 years of Audio Overviews in the first three months.
    \item \textbf{Interactive Mode (Dec 2024):} ``Join the conversation'' — tap the Join button to ask the AI hosts questions mid-overview and receive personalized, real-time responses.
    \item \textbf{Customization (2025):} Choose from four formats — \textit{Deep Dive}, \textit{Brief}, \textit{Critique}, or \textit{Debate}. Adjust tone (formal/humorous) and length. Inject custom instructions to steer the hosts' focus.
    \item \textbf{80+ Languages (Aug 2025):} Expanded from English-only to 80+ languages with full depth parity. Audio Overviews now match the quality of the original English version across all supported languages.
    \item \textbf{Scope Control:} Generate an Audio Overview from all sources in your notebook, or select a specific subset for a focused summary.
  \end{itemize}
\end{frame}

%%%%%%%%%%%%%%%%%%%%%%%%%%%%%%%%%%%%%%%%%%%%%%%%%%%%%%%%%%%%%%%%%%%%%%%%%%%%%%%%%%
\begin{frame}[fragile]\frametitle{Video Overviews, Slide Decks \& Infographics}
  \begin{itemize}
    \item \textbf{Video Overviews (2025):} AI hosts discuss your source material while presenting relevant AI-generated visuals, diagrams, charts, and quotes — creating a narrated slide-style video. Available in six visual styles powered by Google's Nano Banana image generation model and in formats: \textit{Brief} and \textit{Explainer}. Available in 80+ languages.
    \item \textbf{Slide Decks (Nov 2025):} Generate a structured presentation directly from your sources — complete with layouts, visuals, and narrative text. Long-form deck option available on higher tiers. Watermarks appear on Free/Plus/Pro tiers; removed for Ultra subscribers.
    \item \textbf{Infographics (Nov 2025):} Create visually rich summary graphics from your uploaded sources. Powered by Nano Banana image generation. Watermarks removed for Ultra tier.
    \item \textbf{Export:} Slide Decks and Infographics can be exported and shared. Mind Maps export as PNG (for presentations) or JSON (for integrations).
    \item \textbf{Multiple Outputs:} The Studio Panel stores multiple generated outputs of the same type within a single notebook — iterate freely without overwriting previous versions.
  \end{itemize}
\end{frame}

%%%%%%%%%%%%%%%%%%%%%%%%%%%%%%%%%%%%%%%%%%%%%%%%%%%%%%%%%%%%%%%%%%%%%%%%%%%%%%%%%%
\begin{frame}[fragile]\frametitle{Mind Maps, Reports \& Learning Tools}
  \begin{itemize}
    \item \textbf{Mind Maps (Mar 2025):} One click generates an interactive visual map of key concepts and connections extracted from your sources. Nodes expand multi-level deep; click any node to open a targeted AI chat on that topic. Export as PNG or JSON.
    \item \textbf{Reports:} Ask NotebookLM to synthesize findings on a specific question across all sources — delivering a cited, structured briefing document. Useful for literature reviews, executive briefings, and competitive analysis.
    \item \textbf{Data Tables (Dec 2025):} Synthesize scattered data across sources into clean, structured tables — exportable directly to Google Sheets. Use cases: clinical trial comparisons, action items from meeting transcripts, competitor pricing matrices, exam study tables.
    \item \textbf{Flashcards:} Auto-generated from source content; ideal for exam preparation and employee training.
    \item \textbf{Quizzes:} NotebookLM generates multiple-choice and open-ended assessments grounded in your uploaded materials.
    \item \textbf{Study Guides \& Timelines:} Pre-built outputs for academic workflows — chronological event timelines and structured study guides with key terms, questions, and summaries.
  \end{itemize}
\end{frame}

%%%%%%%%%%%%%%%%%%%%%%%%%%%%%%%%%%%%%%%%%%%%%%%%%%%%%%%%%%%%%%%%%%%%%%%%%%%%%%%%%%
\begin{frame}[fragile]\frametitle{Deep Research — Autonomous Source Discovery}
  \begin{itemize}
    \item \textbf{What It Is:} NotebookLM's agentic research mode. Give it a topic and it autonomously creates a research plan, searches hundreds of high-quality web sources, and compiles a comprehensive, citation-backed report — then adds those sources to your notebook.
    \item \textbf{Fast Research vs.\ Deep Research:} \textit{Fast Research} surfaces $\sim$10 relevant documents in under 30 seconds. \textit{Deep Research} runs a broader, multi-step crawl — slower but significantly more comprehensive.
    \item \textbf{Source Discovery Panel:} A dedicated panel within NotebookLM lets you search for sources from the web or from your Google Drive before deciding which to add.
    \item \textbf{Workflow Integration:} Combine external tools — e.g., use Gemini for broad web research, then feed the output into NotebookLM as a grounded source for focused Q\&A and synthesis.
    \item \textbf{Plan Limits:} Free tier allows 10 Deep Research sessions per month; Pro allows up to 20 per day; Ultra allows up to 200 per day.
    \item \textbf{Important Caveat:} If information is not in your sources (uploaded or discovered), NotebookLM will not speculate. It will say so explicitly — a feature, not a bug.
  \end{itemize}
\end{frame}

%%%%%%%%%%%%%%%%%%%%%%%%%%%%%%%%%%%%%%%%%%%%%%%%%%%%%%%%%%%%%%%%%%%%%%%%%%%%%%%%%%
\begin{frame}[fragile]\frametitle{Pricing \& Plans}
  \begin{center}
  \begin{tabular}{|l|c|c|c|c|}
    \hline
    \textbf{Feature} & \textbf{Free} & \textbf{Plus} & \textbf{Pro} & \textbf{Ultra} \\
    \hline
    Price & \$0 & Workspace & \$19.99/mo & \$249.99/mo \\
    Notebooks & 100 & 200 & 500 & 500 \\
    Sources / Notebook & 50 & 100 & 300 & 600 \\
    Chat queries / day & 50 & 2$\times$ & 5$\times$ & 5,000 \\
    Audio Overviews / day & 3 & 2$\times$ & 5$\times$ & 200 \\
    Video Overviews / day & Limited & 2$\times$ & 5$\times$ & 200 \\
    Deep Research / month & 10/mo & 2$\times$ & 20/day & 200/day \\
    Slide Deck watermark & Yes & Yes & Yes & Removed \\
    Model access & Gemini & Higher & Higher & Highest \\
    \hline
  \end{tabular}
  \end{center}
  \vspace{0.3em}
  \begin{itemize}
    \item \textbf{Pro} is bundled with Google AI Pro (formerly Google One AI Premium) — also includes Gemini Advanced, AI in Gmail/Docs, and 2TB storage. US students get 50\% off (\$9.99/mo, first year).
    \item \textbf{Plus} is included with qualifying Google Workspace plans (starting at \$14/user/month).
    \item \textbf{Ultra} also removes watermarks on Slide Decks and Infographics and grants priority feature access.
  \end{itemize}
\end{frame}

%%%%%%%%%%%%%%%%%%%%%%%%%%%%%%%%%%%%%%%%%%%%%%%%%%%%%%%%%%%%%%%%%%%%%%%%%%%%%%%%%%
\begin{frame}[fragile]\frametitle{Enterprise \& Collaboration Features}
  \begin{itemize}
    \item \textbf{NotebookLM Enterprise (GA 2025):} Runs within the customer's own Google Cloud project. Data is not shared externally; processed inside VPC Service Controls and Identity Access Management (IAM). Suitable for regulated industries (legal, healthcare, finance).
    \item \textbf{Enterprise API (GA):} Automate notebook creation and management — e.g., generate a template notebook per new project, auto-insert RFPs, produce initial report skeletons programmatically.
    \item \textbf{Google Agentspace Integration:} Enterprises can deploy NotebookLM as an internal AI agent within Google Agentspace — treating it as a knowledge infrastructure layer, not just a single-user tool.
    \item \textbf{Collaboration:} Share notebooks with teammates as viewers or editors (similar to Google Docs sharing). Publish notebooks publicly. Notebooks can be scoped to only expose the chat/search interface to collaborators, keeping raw sources private.
    \item \textbf{Data Privacy Guarantee:} Uploaded materials, queries, and model responses are \textbf{not} used to train Google's models — on Free, Pro, and Workspace tiers. Confirmed by Google's product documentation.
    \item \textbf{Admin Controls:} Workspace admins can enable/disable NotebookLM at the domain, OU, or group level. Context-Aware Access (CAA) is supported.
  \end{itemize}
\end{frame}

%%%%%%%%%%%%%%%%%%%%%%%%%%%%%%%%%%%%%%%%%%%%%%%%%%%%%%%%%%%%%%%%%%%%%%%%%%%%%%%%%%
\begin{frame}[fragile]\frametitle{Key Use Cases}
  \begin{itemize}
    \item \textbf{Academic Research:} Accelerate literature reviews — upload papers, generate mind maps of themes, query across all PDFs simultaneously with cited answers, produce timelines and study guides. Every answer links back to the exact passage.
    \item \textbf{Enterprise Knowledge Management:} Upload compliance manuals, product docs, HR policies. New hires query the notebook instead of hunting through documents. Onboarding time reduced significantly.
    \item \textbf{Sales \& Presales:} Digest RFPs, specs, and past proposals; generate tailored meeting plans; produce audio/video briefings to prepare for customer conversations.
    \item \textbf{Executive Intelligence:} Upload financial reports, market analyses, and strategy documents; generate executive summaries, extract metrics, and identify strategic implications.
    \item \textbf{Content Production:} Turn research bundles into audio guides, infographics, slide decks, and video explainers — without manual formatting work.
    \item \textbf{IT \& Engineering:} Upload production logs and post-mortem documents; generate root cause timelines and mind maps of system dependencies; query issues across hundreds of pages of logs instantly.
    \item \textbf{Education:} Generate quizzes, flashcards, and study guides from course materials; create audio overviews for on-the-go learning.
  \end{itemize}
\end{frame}

%%%%%%%%%%%%%%%%%%%%%%%%%%%%%%%%%%%%%%%%%%%%%%%%%%%%%%%%%%%%%%%%%%%%%%%%%%%%%%%%%%
\begin{frame}[fragile]\frametitle{Limitations \& Honest Caveats}
  \begin{itemize}
    \item \textbf{Cloud-Only:} No offline mode on any plan. All processing requires an active internet connection and happens on Google's infrastructure.
    \item \textbf{Single AI Provider:} Only Google Gemini models are available. You cannot bring your own API keys, switch to OpenAI or Claude, or run local models.
    \item \textbf{Static Sources:} Documents do not update automatically. When a source file changes, you must manually remove the old version and re-upload the new one.
    \item \textbf{No Chat History:} NotebookLM does not retain conversation history between sessions. Save important exchanges as notes and convert them to sources if needed.
    \item \textbf{Source-Bound Answers:} If information is not in your uploaded sources, NotebookLM will not answer from general knowledge. This is a deliberate design choice for accuracy, but limits open-ended queries.
    \item \textbf{Quality Degrades Near Source Limit:} Answer accuracy declines as you approach your plan's source cap — the model retrieves from an increasingly large pool, reducing precision.
    \item \textbf{Safety Filters:} Queries involving violence, sexuality, or obscenity in sources (even in historical contexts) may trigger safety flags and prevent responses.
    \item \textbf{Google Ecosystem Lock-In:} Deep integration with Drive, Docs, Sheets, and Slides is a strength for Google Workspace users but a constraint for Microsoft 365-centric organizations.
  \end{itemize}
\end{frame}

%%%%%%%%%%%%%%%%%%%%%%%%%%%%%%%%%%%%%%%%%%%%%%%%%%%%%%%%%%%%%%%%%%%%%%%%%%%%%%%%%%
\begin{frame}[fragile]\frametitle{Demo 1 — Research Synthesis in Minutes}
  \begin{itemize}
    \item \textbf{Goal:} Show core source grounding, cited Q\&A, and Audio Overview generation.
    \item \textbf{Setup:} Create a new notebook. Add 5--10 PDFs on a research topic (e.g., AI agent architectures) plus two or three YouTube video URLs and a Google Doc of notes.
    \item \textbf{Step 1 — Explore:} Click ``Generate'' on the Sources panel. NotebookLM suggests relevant questions based on the uploaded material — select one.
    \item \textbf{Step 2 — Q\&A with Citations:} Ask:
\begin{lstlisting}
> What are the key differences in memory
  architecture between the papers?
\end{lstlisting}
    Show inline citations — click one to jump to the exact passage in the source PDF.
    \item \textbf{Step 3 — Audio Overview:} Studio panel → Audio Overview → choose ``Deep Dive'' format. Play the generated podcast. Demonstrate the ``Join'' button to ask a live question to the AI hosts mid-overview.
    \item \textbf{Key Message:} From raw documents to a listenable, cited research briefing in under two minutes — with zero hallucination risk.
  \end{itemize}
\end{frame}

%%%%%%%%%%%%%%%%%%%%%%%%%%%%%%%%%%%%%%%%%%%%%%%%%%%%%%%%%%%%%%%%%%%%%%%%%%%%%%%%%%
\begin{frame}[fragile]\frametitle{Demo 2 — Mind Map \& Data Table}
  \begin{itemize}
    \item \textbf{Goal:} Show visual synthesis and structured data extraction.
    \item \textbf{Setup:} Notebook with 10--15 market research reports and competitor analysis PDFs.
    \item \textbf{Step 1 — Mind Map:} Studio panel → Mind Map. Show the auto-generated concept graph. Click a top-level node (e.g., ``Pricing Strategies'') to expand sub-nodes. Click a sub-node to open a targeted chat: the AI answers questions scoped to that concept cluster.
    \item \textbf{Step 2 — Data Table:} In the Chat panel, type:
\begin{lstlisting}
> Create a table comparing each competitor's
  pricing tier, feature set, and target
  customer segment from the sources.
\end{lstlisting}
    Show the structured table output. Click ``Export to Google Sheets'' to transfer it directly.
    \item \textbf{What to Highlight:} The mind map and table are generated from the same source set — zero manual data entry, full citation traceability.
    \item \textbf{Key Message:} NotebookLM turns a folder of reports into an interactive knowledge graph and a structured data asset simultaneously.
  \end{itemize}
\end{frame}

%%%%%%%%%%%%%%%%%%%%%%%%%%%%%%%%%%%%%%%%%%%%%%%%%%%%%%%%%%%%%%%%%%%%%%%%%%%%%%%%%%
\begin{frame}[fragile]\frametitle{Demo 3 — Slide Deck \& Video Overview}
  \begin{itemize}
    \item \textbf{Goal:} Show multimodal output generation for presentations and async communication.
    \item \textbf{Setup:} Notebook with a Deep Research report, three supporting PDFs, and a meeting transcript.
    \item \textbf{Step 1 — Slide Deck:} Studio panel → Slide Deck → choose ``Long'' format (Pro/Ultra). Show the generated presentation with layouts, AI-generated visuals, and source-derived text. Note the watermark (removed on Ultra).
    \item \textbf{Step 2 — Video Overview:} Studio panel → Video Overview → choose visual style (e.g., ``Explainer''). Play the AI-narrated video with slides, diagrams, and quotes drawn from sources.
    \item \textbf{Step 3 — Output Language:} Demonstrate switching the output language selector to Japanese — regenerate the briefing document in Japanese from the same English sources.
    \item \textbf{What to Show:} All three outputs (Slide Deck, Video, Translated Briefing) generated from a single source set in under five minutes.
    \item \textbf{Key Message:} NotebookLM is not just a research tool — it is a multimodal content production engine grounded in verified sources.
  \end{itemize}
\end{frame}

%%%%%%%%%%%%%%%%%%%%%%%%%%%%%%%%%%%%%%%%%%%%%%%%%%%%%%%%%%%%%%%%%%%%%%%%%%%%%%%%%%
\begin{frame}[fragile]\frametitle{Demo 4 — Enterprise: Onboarding \& Deep Research}
  \begin{itemize}
    \item \textbf{Goal:} Show collaborative notebooks and autonomous source discovery for enterprise workflows.
    \item \textbf{Setup:} A shared team notebook with product specs, onboarding manuals, and policy documents (Google Workspace account — data stays within org trust boundary).
    \item \textbf{Step 1 — Collaborative Q\&A:} Share the notebook with a colleague as a viewer (chat-only access; sources remain private). Show the colleague asking a policy question and getting a cited answer without seeing the underlying documents.
    \item \textbf{Step 2 — Deep Research:} In a separate research notebook, use the Source Discovery panel → Deep Research. Enter:
\begin{lstlisting}
> Competitive landscape for enterprise AI
  research assistants, Q4 2025
\end{lstlisting}
    Show NotebookLM building a bibliography autonomously, then generate a cited briefing report.
    \item \textbf{Step 3 — Admin Controls:} Briefly show the Google Workspace Admin console toggle for enabling/disabling NotebookLM at the domain level.
    \item \textbf{Key Message:} Enterprise-grade access control, data governance, and collaborative research — all within the Google Workspace security boundary.
  \end{itemize}
\end{frame}

%%%%%%%%%%%%%%%%%%%%%%%%%%%%%%%%%%%%%%%%%%%%%%%%%%%%%%%%%%%%%%%%%%%%%%%%%%%%%%%%%%
\begin{frame}[fragile]\frametitle{NotebookLM vs.\ Competing Tools}
  \begin{center}
  \begin{tabular}{|l|c|c|c|}
    \hline
    \textbf{Feature} & \textbf{NotebookLM} & \textbf{ChatGPT Projects} & \textbf{Claude} \\
    \hline
    Source Grounding & Yes (RAG, cited) & Partial & Partial \\
    Inline Citations & Yes (passage-level) & No & No \\
    Audio Podcast Output & Yes & No & No \\
    Video Overview & Yes & No & No \\
    Mind Map & Yes & No & No \\
    Slide Deck Generation & Yes & No & No \\
    Deep Research Agent & Yes & Yes (web) & Yes (web) \\
    Model Choice & Gemini only & GPT only & Claude only \\
    Offline / Local & No & No & No \\
    Free Tier & Yes & Limited & Yes \\
    Data Not Used for Training & Yes & Opt-out & Yes \\
    \hline
  \end{tabular}
  \end{center}
\end{frame}

%%%%%%%%%%%%%%%%%%%%%%%%%%%%%%%%%%%%%%%%%%%%%%%%%%%%%%%%%%%%%%%%%%%%%%%%%%%%%%%%%%
\begin{frame}[fragile]\frametitle{Summary: Why NotebookLM Matters}
  \begin{itemize}
    \item \textbf{Paradigm Shift:} From ``AI that knows everything'' to ``AI that knows \textit{your} documents'' — source-grounded, hallucination-resistant, and fully cited.
    \item \textbf{Multimodal Knowledge Engine:} One notebook, many outputs — audio podcasts, narrated videos, slide decks, infographics, mind maps, data tables, quizzes, and reports — all from the same source set.
    \item \textbf{Accessible \& Free:} Core features available to anyone with a Google account. Powerful enough for serious research; simple enough for first-time AI users.
    \item \textbf{Scales to Enterprise:} VPC-SC, IAM, admin controls, and an enterprise API make it deployable in regulated environments where data cannot leave the organization.
    \item \textbf{Viral by Design:} Audio Overviews became a cultural moment — 350+ years of listening generated in the first months. Spotify integration. NotebookLM proved AI can be genuinely delightful.
    \item \textbf{Trajectory:} 2023--24 solved the input problem (RAG). 2025 solved the format problem (multimodal outputs). The roadmap points toward persistent memory, agentic orchestration (Google Agentspace), and ambient AI through Project Astra / XR devices.
    \item \textbf{The Bigger Picture:} NotebookLM is Google's answer to ``What if AI only ever spoke from documents you trust?''
  \end{itemize}
\end{frame}

%%%%%%%%%%%%%%%%%%%%%%%%%%%%%%%%%%%%%%%%%%%%%%%%%%%%%%%%%%%%%%%%%%%%%%%%%%%%%%%%%%
\begin{frame}[fragile]\frametitle{Resources \& Links}
  \begin{itemize}
    \item \textbf{Product:} \href{https://notebooklm.google.com}{\texttt{notebooklm.google.com}}
    \item \textbf{Plans \& Pricing:} \href{https://notebooklm.google/plans}{\texttt{notebooklm.google/plans}}
    \item \textbf{Help Center — FAQ:} \href{https://support.google.com/notebooklm/answer/16269187}{\texttt{support.google.com/notebooklm → FAQ}}
    \item \textbf{Help Center — Upgrade:} \href{https://support.google.com/notebooklm/answer/16213268}{\texttt{support.google.com/notebooklm → Upgrade}}
    \item \textbf{Google Workspace:} \href{https://workspace.google.com/products/notebooklm/}{\texttt{workspace.google.com/products/notebooklm}}
    \item \textbf{Workspace Updates Blog:} \href{https://workspaceupdates.googleblog.com}{\texttt{workspaceupdates.googleblog.com}} (search: NotebookLM)
    \item \textbf{Google Labs Blog:} \href{https://blog.google/innovation-and-ai/models-and-research/google-labs/}{\texttt{blog.google → Google Labs}} (latest release notes)
    \item \textbf{Mobile App:} Available on iOS and Android — core features including source upload via one-tap share, Audio Overview playback, and chat.
  \end{itemize}
\end{frame}
